% --- Archon physics formalization source begin ---
% archon:physics
% archon:covers PhyXMiniProblems/problem_phyx_mini_0630.lean
% archon:source-report reports/phyx_mini/problem_phyx_mini_0630.source.json

\chapter{Physics problem phyx\_mini\_0630}
\label{ch:PhyXMiniProblems_problem_phyx_mini_0630}

\paragraph{Problem source.}
\#\# Physical scenario

A \$p\$-\$n\$ junction includes \$p\$-type silicon, with donor atom density \$N\_D\$, adjacent to \$n\$-type silicon, with acceptor atom density \$N\_A\$. Near the junction, free electrons from the \$n\$ side have diffused into the \$p\$ side while free holes from the \$p\$ side have diffused into the \$n\$ side, leaving a “depletion region” of width \$W\$ where there are no free charge carriers. The width of the depletion region on the \$n\$ side is \$a\_n\$, and the width of the depletion region on the \$p\$ side is \$a\_p\$. Each depleted region has a constant charge density with magnitude equal to the corresponding donor atom density. Assume there is no net charge outside the depleted regions. Define an \$x\$-axis pointing from the \$p\$ side toward the \$n\$ side with the origin at the center of the junction. An electric field \$\textbackslash{}vec\{E\} = -E\textbackslash{}hat\{\textbackslash{}imath\}\$ has developed in the depletion region, stabilizing the diffusion.

\#\# Question

The \$p\$ side is doped with boron atoms with density \$N\_A = 1.00 \textbackslash{}times 10\textasciicircum{}\{16\}\textbackslash{} cm\textasciicircum{}\{-3\}\$. The \$n\$ side is doped with arsenic atoms with density \$N\_D = 5.00 \textbackslash{}times 10\textasciicircum{}\{16\}\textbackslash{} cm\textasciicircum{}\{-3\}\$. The \$n\$ side depletion depth is \$a\_n = 55.0\textbackslash{} nm\$. What is the \$p\$ side depletion depth \$a\_p\$?

\#\# Answer choices

- A: A: 275nm

- B: B: 399nm

- C: C: 250nm

- D: D: 0.350nm

\#\# Figure caption (auxiliary; use the image as primary evidence)

The image illustrates a segment of a semiconductor device, likely a diode, showing a cross-sectional view of a p-n junction. It is divided into three main regions:

1. **Left Region (p-type):** This section is shaded in a light peach color and labeled as "p". It is the p-type semiconductor, enriched with positive charge carriers (holes).

2. **Middle Region (Depletion Layer):** Located between the p-type and n-type regions, this section is represented with two vertical areas:
   - The left part is within the p-type region and has dashed lines, indicating the presence of negative charges.
   - The right part is within the n-type region and has plus symbols, indicating the presence of positive charges.
   - This entire region is labeled with "0" at its left boundary, and its width is denoted as "w". The boundaries are marked as \textbackslash{}(-a\_p\textbackslash{}) and \textbackslash{}(a\_n\textbackslash{}).

3. **Right Region (n-type):** This section is shaded in light blue and labeled as "n". It is the n-type semiconductor, enriched with negative charge carriers (electrons).

Arrows enter each end of the diagram, signifying the direction of current flow. The entire illustration represents the junction formed between these semiconductor types, with the depletion region indicating charge separation at equilibrium.

\#\# Dataset metadata

Quantum Phenomena; Physical Model Grounding Reasoning, Multi-Formula Reasoning

\paragraph{Recorded answer/context.}
A: A: 275nm

\paragraph{Figure/image path.}
/root/proposal\_for\_physic/science-mango/phyx\_mini\_run/phyx\_data/test\_image/630.png

\paragraph{Formalization target.}
create a compiling Lean file with sorry bodies at `PhyXMiniProblems/problem\_phyx\_mini\_0630.lean`. The Lean declarations must preserve the physical quantities, dimensions or dimensional roles, figure labels, governing-law hypotheses, and final relation expressed by this problem.
Use Mathlib/Physlib names found through LeanExplore where available. If a physics API is missing, introduce faithful local abstractions rather than scalar placeholder aliases.

\begin{theorem}[Physics formalization target]
\label{thm:physics:phyx_mini_0630:target}
\lean{PhyXMiniProblems.ProblemPhyXMini0630.p_side_depletion_depth_is_275_nanometers}
\uses{def:physics:phyx-mini-0630:phyxminiproblems-problemphyxmini0630-lengthinnanometers, def:physics:phyx-mini-0630:phyxminiproblems-problemphyxmini0630-pnjunctionsetup, def:physics:phyx-mini-0630:phyxminiproblems-problemphyxmini0630-matchesproblemdata, def:physics:phyx-mini-0630:phyxminiproblems-problemphyxmini0630-matchesprimaryfigure, def:physics:phyx-mini-0630:phyxminiproblems-problemphyxmini0630-hasphysicalparameters, def:physics:phyx-mini-0630:phyxminiproblems-problemphyxmini0630-satisfiesdepletionregionphysics, def:physics:phyx-mini-0630:phyxminiproblems-problemphyxmini0630-matchesanswerchoice}
The assigned autoformalize agent should translate this physics problem into Lean declarations in the covered file, with theorem and lemma proof bodies written as `by sorry`.
\end{theorem}
\begin{proof}
This is an autoformalization task, not a proof task. Produce statements that can later be proved by the physics prover without weakening the physical model.
\end{proof}

% --- Archon declaration topology begin ---
\section{Declaration topology}
\begin{definition}[inverse Volume Dimension]
  \label{def:physics:phyx-mini-0630:phyxminiproblems-problemphyxmini0630-inversevolumedimension}
  \lean{PhyXMiniProblems.ProblemPhyXMini0630.inverseVolumeDimension}
  The inverse-volume dimension carried by a number density
\end{definition}
\begin{proof}
  This is the declared model definition of inverse Volume Dimension.
\end{proof}
\begin{definition}[charge Density Dimension]
  \label{def:physics:phyx-mini-0630:phyxminiproblems-problemphyxmini0630-chargedensitydimension}
  \lean{PhyXMiniProblems.ProblemPhyXMini0630.chargeDensityDimension}
  \uses{def:physics:phyx-mini-0630:phyxminiproblems-problemphyxmini0630-inversevolumedimension}
  The dimension of electric charge per unit volume
\end{definition}
\begin{proof}
  This is the declared model definition of charge Density Dimension.
\end{proof}
\begin{definition}[energy Dimension]
  \label{def:physics:phyx-mini-0630:phyxminiproblems-problemphyxmini0630-energydimension}
  \lean{PhyXMiniProblems.ProblemPhyXMini0630.energyDimension}
  The dimension of energy
\end{definition}
\begin{proof}
  This is the declared model definition of energy Dimension.
\end{proof}
\begin{definition}[electric Field Dimension]
  \label{def:physics:phyx-mini-0630:phyxminiproblems-problemphyxmini0630-electricfielddimension}
  \lean{PhyXMiniProblems.ProblemPhyXMini0630.electricFieldDimension}
  \uses{def:physics:phyx-mini-0630:phyxminiproblems-problemphyxmini0630-energydimension}
  The dimension of one Cartesian component of an electric field
\end{definition}
\begin{proof}
  This is the declared model definition of electric Field Dimension.
\end{proof}
\begin{definition}[Length Quantity]
  \label{def:physics:phyx-mini-0630:phyxminiproblems-problemphyxmini0630-lengthquantity}
  \lean{PhyXMiniProblems.ProblemPhyXMini0630.LengthQuantity}
  A nonnegative physical length, used for depletion depths and widths
\end{definition}
\begin{proof}
  This is the declared model definition of Length Quantity.
\end{proof}
\begin{definition}[Signed Axial Position Quantity]
  \label{def:physics:phyx-mini-0630:phyxminiproblems-problemphyxmini0630-signedaxialpositionquantity}
  \lean{PhyXMiniProblems.ProblemPhyXMini0630.SignedAxialPositionQuantity}
  A signed physical coordinate on the junction axis
\end{definition}
\begin{proof}
  This is the declared model definition of Signed Axial Position Quantity.
\end{proof}
\begin{definition}[Number Density Quantity]
  \label{def:physics:phyx-mini-0630:phyxminiproblems-problemphyxmini0630-numberdensityquantity}
  \lean{PhyXMiniProblems.ProblemPhyXMini0630.NumberDensityQuantity}
  \uses{def:physics:phyx-mini-0630:phyxminiproblems-problemphyxmini0630-inversevolumedimension}
  A nonnegative atom or mobile-carrier number density
\end{definition}
\begin{proof}
  This is the declared model definition of Number Density Quantity.
\end{proof}
\begin{definition}[Charge Magnitude Quantity]
  \label{def:physics:phyx-mini-0630:phyxminiproblems-problemphyxmini0630-chargemagnitudequantity}
  \lean{PhyXMiniProblems.ProblemPhyXMini0630.ChargeMagnitudeQuantity}
  A nonnegative magnitude of electric charge
\end{definition}
\begin{proof}
  This is the declared model definition of Charge Magnitude Quantity.
\end{proof}
\begin{definition}[Charge Density Quantity]
  \label{def:physics:phyx-mini-0630:phyxminiproblems-problemphyxmini0630-chargedensityquantity}
  \lean{PhyXMiniProblems.ProblemPhyXMini0630.ChargeDensityQuantity}
  \uses{def:physics:phyx-mini-0630:phyxminiproblems-problemphyxmini0630-chargedensitydimension}
  A signed volume charge density
\end{definition}
\begin{proof}
  This is the declared model definition of Charge Density Quantity.
\end{proof}
\begin{definition}[Axial Electric Field Quantity]
  \label{def:physics:phyx-mini-0630:phyxminiproblems-problemphyxmini0630-axialelectricfieldquantity}
  \lean{PhyXMiniProblems.ProblemPhyXMini0630.AxialElectricFieldQuantity}
  \uses{def:physics:phyx-mini-0630:phyxminiproblems-problemphyxmini0630-electricfielddimension}
  A signed x-component of the electric field
\end{definition}
\begin{proof}
  This is the declared model definition of Axial Electric Field Quantity.
\end{proof}
\begin{definition}[length Readout]
  \label{def:physics:phyx-mini-0630:phyxminiproblems-problemphyxmini0630-lengthreadout}
  \lean{PhyXMiniProblems.ProblemPhyXMini0630.lengthReadout}
  \uses{def:physics:phyx-mini-0630:phyxminiproblems-problemphyxmini0630-lengthquantity}
  Read a nonnegative physical length in the selected length unit
\end{definition}
\begin{proof}
  This is the declared model definition of length Readout.
\end{proof}
\begin{definition}[signed Position Readout]
  \label{def:physics:phyx-mini-0630:phyxminiproblems-problemphyxmini0630-signedpositionreadout}
  \lean{PhyXMiniProblems.ProblemPhyXMini0630.signedPositionReadout}
  \uses{def:physics:phyx-mini-0630:phyxminiproblems-problemphyxmini0630-signedaxialpositionquantity}
  Read a signed axial coordinate in the selected length unit
\end{definition}
\begin{proof}
  This is the declared model definition of signed Position Readout.
\end{proof}
\begin{definition}[number Density Readout]
  \label{def:physics:phyx-mini-0630:phyxminiproblems-problemphyxmini0630-numberdensityreadout}
  \lean{PhyXMiniProblems.ProblemPhyXMini0630.numberDensityReadout}
  \uses{def:physics:phyx-mini-0630:phyxminiproblems-problemphyxmini0630-numberdensityquantity}
  Read a physical number density in inverse cubes of the selected unit
\end{definition}
\begin{proof}
  This is the declared model definition of number Density Readout.
\end{proof}
\begin{definition}[charge Density Readout]
  \label{def:physics:phyx-mini-0630:phyxminiproblems-problemphyxmini0630-chargedensityreadout}
  \lean{PhyXMiniProblems.ProblemPhyXMini0630.chargeDensityReadout}
  \uses{def:physics:phyx-mini-0630:phyxminiproblems-problemphyxmini0630-chargedensityquantity}
  Read a charge density in coulombs per cube of the selected length unit
\end{definition}
\begin{proof}
  This is the declared model definition of charge Density Readout.
\end{proof}
\begin{definition}[charge Magnitude In Coulombs]
  \label{def:physics:phyx-mini-0630:phyxminiproblems-problemphyxmini0630-chargemagnitudeincoulombs}
  \lean{PhyXMiniProblems.ProblemPhyXMini0630.chargeMagnitudeInCoulombs}
  \uses{def:physics:phyx-mini-0630:phyxminiproblems-problemphyxmini0630-chargemagnitudequantity}
  Coulomb readout of a nonnegative physical charge magnitude
\end{definition}
\begin{proof}
  This is the declared model definition of charge Magnitude In Coulombs.
\end{proof}
\begin{definition}[axial Electric Field In Volts Per Meter]
  \label{def:physics:phyx-mini-0630:phyxminiproblems-problemphyxmini0630-axialelectricfieldinvoltspermeter}
  \lean{PhyXMiniProblems.ProblemPhyXMini0630.axialElectricFieldInVoltsPerMeter}
  \uses{def:physics:phyx-mini-0630:phyxminiproblems-problemphyxmini0630-axialelectricfieldquantity}
  Coherent-SI readout of a signed electric-field component
\end{definition}
\begin{proof}
  This is the declared model definition of axial Electric Field In Volts Per Meter.
\end{proof}
\begin{definition}[length In Nanometers]
  \label{def:physics:phyx-mini-0630:phyxminiproblems-problemphyxmini0630-lengthinnanometers}
  \lean{PhyXMiniProblems.ProblemPhyXMini0630.lengthInNanometers}
  \uses{def:physics:phyx-mini-0630:phyxminiproblems-problemphyxmini0630-lengthquantity, def:physics:phyx-mini-0630:phyxminiproblems-problemphyxmini0630-lengthreadout}
  Nanometre readout used for the supplied depth and answer choices
\end{definition}
\begin{proof}
  This is the declared model definition of length In Nanometers.
\end{proof}
\begin{definition}[signed Position In Nanometers]
  \label{def:physics:phyx-mini-0630:phyxminiproblems-problemphyxmini0630-signedpositioninnanometers}
  \lean{PhyXMiniProblems.ProblemPhyXMini0630.signedPositionInNanometers}
  \uses{def:physics:phyx-mini-0630:phyxminiproblems-problemphyxmini0630-signedaxialpositionquantity, def:physics:phyx-mini-0630:phyxminiproblems-problemphyxmini0630-signedpositionreadout}
  Nanometre readout of a signed axial coordinate
\end{definition}
\begin{proof}
  This is the declared model definition of signed Position In Nanometers.
\end{proof}
\begin{definition}[length In Centimeters]
  \label{def:physics:phyx-mini-0630:phyxminiproblems-problemphyxmini0630-lengthincentimeters}
  \lean{PhyXMiniProblems.ProblemPhyXMini0630.lengthInCentimeters}
  \uses{def:physics:phyx-mini-0630:phyxminiproblems-problemphyxmini0630-lengthquantity, def:physics:phyx-mini-0630:phyxminiproblems-problemphyxmini0630-lengthreadout}
  Centimetre readout used in the areal charge-neutrality law
\end{definition}
\begin{proof}
  This is the declared model definition of length In Centimeters.
\end{proof}
\begin{definition}[number Density Per Cubic Centimeter]
  \label{def:physics:phyx-mini-0630:phyxminiproblems-problemphyxmini0630-numberdensitypercubiccentimeter}
  \lean{PhyXMiniProblems.ProblemPhyXMini0630.numberDensityPerCubicCentimeter}
  \uses{def:physics:phyx-mini-0630:phyxminiproblems-problemphyxmini0630-numberdensityquantity, def:physics:phyx-mini-0630:phyxminiproblems-problemphyxmini0630-numberdensityreadout}
  Atom-density readout in the cm\textasciicircum{}-3 units stated in the problem
\end{definition}
\begin{proof}
  This is the declared model definition of number Density Per Cubic Centimeter.
\end{proof}
\begin{definition}[charge Density In Coulombs Per Cubic Centimeter]
  \label{def:physics:phyx-mini-0630:phyxminiproblems-problemphyxmini0630-chargedensityincoulombspercubiccentimeter}
  \lean{PhyXMiniProblems.ProblemPhyXMini0630.chargeDensityInCoulombsPerCubicCentimeter}
  \uses{def:physics:phyx-mini-0630:phyxminiproblems-problemphyxmini0630-chargedensityquantity, def:physics:phyx-mini-0630:phyxminiproblems-problemphyxmini0630-chargedensityreadout}
  Charge-density readout in coulombs per cubic centimetre
\end{definition}
\begin{proof}
  This is the declared model definition of charge Density In Coulombs Per Cubic Centimeter.
\end{proof}
\begin{definition}[Semiconductor Side]
  \label{def:physics:phyx-mini-0630:phyxminiproblems-problemphyxmini0630-semiconductorside}
  \lean{PhyXMiniProblems.ProblemPhyXMini0630.SemiconductorSide}
  The two sides of the junction
\end{definition}
\begin{proof}
  This is the declared model definition of Semiconductor Side.
\end{proof}
\begin{definition}[Conductivity Type]
  \label{def:physics:phyx-mini-0630:phyxminiproblems-problemphyxmini0630-conductivitytype}
  \lean{PhyXMiniProblems.ProblemPhyXMini0630.ConductivityType}
  Conductivity types named in the source
\end{definition}
\begin{proof}
  This is the declared model definition of Conductivity Type.
\end{proof}
\begin{definition}[Semiconductor Material]
  \label{def:physics:phyx-mini-0630:phyxminiproblems-problemphyxmini0630-semiconductormaterial}
  \lean{PhyXMiniProblems.ProblemPhyXMini0630.SemiconductorMaterial}
  Semiconductor material used on both sides
\end{definition}
\begin{proof}
  This is the declared model definition of Semiconductor Material.
\end{proof}
\begin{definition}[Dopant Species]
  \label{def:physics:phyx-mini-0630:phyxminiproblems-problemphyxmini0630-dopantspecies}
  \lean{PhyXMiniProblems.ProblemPhyXMini0630.DopantSpecies}
  Dopant species supplied in the numerical question
\end{definition}
\begin{proof}
  This is the declared model definition of Dopant Species.
\end{proof}
\begin{definition}[Mobile Carrier]
  \label{def:physics:phyx-mini-0630:phyxminiproblems-problemphyxmini0630-mobilecarrier}
  \lean{PhyXMiniProblems.ProblemPhyXMini0630.MobileCarrier}
  Mobile charge carriers excluded from the depletion region
\end{definition}
\begin{proof}
  This is the declared model definition of Mobile Carrier.
\end{proof}
\begin{definition}[Fixed Charge Sign]
  \label{def:physics:phyx-mini-0630:phyxminiproblems-problemphyxmini0630-fixedchargesign}
  \lean{PhyXMiniProblems.ProblemPhyXMini0630.FixedChargeSign}
  Signs of the fixed ions drawn in the two depleted regions
\end{definition}
\begin{proof}
  This is the declared model definition of Fixed Charge Sign.
\end{proof}
\begin{definition}[Axial Direction]
  \label{def:physics:phyx-mini-0630:phyxminiproblems-problemphyxmini0630-axialdirection}
  \lean{PhyXMiniProblems.ProblemPhyXMini0630.AxialDirection}
  The positive x-axis points from the p side toward the n side
\end{definition}
\begin{proof}
  This is the declared model definition of Axial Direction.
\end{proof}
\begin{definition}[Junction Landmark]
  \label{def:physics:phyx-mini-0630:phyxminiproblems-problemphyxmini0630-junctionlandmark}
  \lean{PhyXMiniProblems.ProblemPhyXMini0630.JunctionLandmark}
  The three physical landmarks labelled along the bottom of the figure
\end{definition}
\begin{proof}
  This is the declared model definition of Junction Landmark.
\end{proof}
\begin{definition}[PNJunction Figure]
  \label{def:physics:phyx-mini-0630:phyxminiproblems-problemphyxmini0630-pnjunctionfigure}
  \lean{PhyXMiniProblems.ProblemPhyXMini0630.PNJunctionFigure}
  \uses{def:physics:phyx-mini-0630:phyxminiproblems-problemphyxmini0630-semiconductorside, def:physics:phyx-mini-0630:phyxminiproblems-problemphyxmini0630-fixedchargesign, def:physics:phyx-mini-0630:phyxminiproblems-problemphyxmini0630-junctionlandmark}
  Qualitative evidence in image 630.png. The horizontal coordinates record only visual ordering; the image does not provide a quantitative spatial scale
\end{definition}
\begin{proof}
  This is the declared model definition of PNJunction Figure.
\end{proof}
\begin{definition}[PNJunction Setup]
  \label{def:physics:phyx-mini-0630:phyxminiproblems-problemphyxmini0630-pnjunctionsetup}
  \lean{PhyXMiniProblems.ProblemPhyXMini0630.PNJunctionSetup}
  \uses{def:physics:phyx-mini-0630:phyxminiproblems-problemphyxmini0630-lengthquantity, def:physics:phyx-mini-0630:phyxminiproblems-problemphyxmini0630-signedaxialpositionquantity, def:physics:phyx-mini-0630:phyxminiproblems-problemphyxmini0630-numberdensityquantity, def:physics:phyx-mini-0630:phyxminiproblems-problemphyxmini0630-chargemagnitudequantity, def:physics:phyx-mini-0630:phyxminiproblems-problemphyxmini0630-chargedensityquantity, def:physics:phyx-mini-0630:phyxminiproblems-problemphyxmini0630-axialelectricfieldquantity, def:physics:phyx-mini-0630:phyxminiproblems-problemphyxmini0630-semiconductorside, def:physics:phyx-mini-0630:phyxminiproblems-problemphyxmini0630-conductivitytype, def:physics:phyx-mini-0630:phyxminiproblems-problemphyxmini0630-semiconductormaterial, def:physics:phyx-mini-0630:phyxminiproblems-problemphyxmini0630-dopantspecies, def:physics:phyx-mini-0630:phyxminiproblems-problemphyxmini0630-mobilecarrier, def:physics:phyx-mini-0630:phyxminiproblems-problemphyxmini0630-axialdirection, def:physics:phyx-mini-0630:phyxminiproblems-problemphyxmini0630-junctionlandmark, def:physics:phyx-mini-0630:phyxminiproblems-problemphyxmini0630-pnjunctionfigure}
  Independent physical quantities for the junction. In particular, pSideDepletionDepth is not defined using an answer value. The charge and field profiles are physical fields over dimensionful axial positions
\end{definition}
\begin{proof}
  This is the declared model definition of PNJunction Setup.
\end{proof}
\begin{definition}[In PSide Depletion]
  \label{def:physics:phyx-mini-0630:phyxminiproblems-problemphyxmini0630-inpsidedepletion}
  \lean{PhyXMiniProblems.ProblemPhyXMini0630.InPSideDepletion}
  \uses{def:physics:phyx-mini-0630:phyxminiproblems-problemphyxmini0630-signedaxialpositionquantity, def:physics:phyx-mini-0630:phyxminiproblems-problemphyxmini0630-lengthinnanometers, def:physics:phyx-mini-0630:phyxminiproblems-problemphyxmini0630-signedpositioninnanometers, def:physics:phyx-mini-0630:phyxminiproblems-problemphyxmini0630-pnjunctionsetup}
  A position is in the depleted p side, from -a\_p up to the junction
\end{definition}
\begin{proof}
  This is the declared model definition of In PSide Depletion.
\end{proof}
\begin{definition}[In NSide Depletion]
  \label{def:physics:phyx-mini-0630:phyxminiproblems-problemphyxmini0630-innsidedepletion}
  \lean{PhyXMiniProblems.ProblemPhyXMini0630.InNSideDepletion}
  \uses{def:physics:phyx-mini-0630:phyxminiproblems-problemphyxmini0630-signedaxialpositionquantity, def:physics:phyx-mini-0630:phyxminiproblems-problemphyxmini0630-lengthinnanometers, def:physics:phyx-mini-0630:phyxminiproblems-problemphyxmini0630-signedpositioninnanometers, def:physics:phyx-mini-0630:phyxminiproblems-problemphyxmini0630-pnjunctionsetup}
  A position is in the depleted n side, from the junction through a\_n
\end{definition}
\begin{proof}
  This is the declared model definition of In NSide Depletion.
\end{proof}
\begin{definition}[In Depletion Region]
  \label{def:physics:phyx-mini-0630:phyxminiproblems-problemphyxmini0630-indepletionregion}
  \lean{PhyXMiniProblems.ProblemPhyXMini0630.InDepletionRegion}
  \uses{def:physics:phyx-mini-0630:phyxminiproblems-problemphyxmini0630-signedaxialpositionquantity, def:physics:phyx-mini-0630:phyxminiproblems-problemphyxmini0630-pnjunctionsetup, def:physics:phyx-mini-0630:phyxminiproblems-problemphyxmini0630-inpsidedepletion, def:physics:phyx-mini-0630:phyxminiproblems-problemphyxmini0630-innsidedepletion}
  The complete depletion region is the union of its p and n parts
\end{definition}
\begin{proof}
  This is the declared model definition of In Depletion Region.
\end{proof}
\begin{definition}[Outside Depletion Region]
  \label{def:physics:phyx-mini-0630:phyxminiproblems-problemphyxmini0630-outsidedepletionregion}
  \lean{PhyXMiniProblems.ProblemPhyXMini0630.OutsideDepletionRegion}
  \uses{def:physics:phyx-mini-0630:phyxminiproblems-problemphyxmini0630-signedaxialpositionquantity, def:physics:phyx-mini-0630:phyxminiproblems-problemphyxmini0630-lengthinnanometers, def:physics:phyx-mini-0630:phyxminiproblems-problemphyxmini0630-signedpositioninnanometers, def:physics:phyx-mini-0630:phyxminiproblems-problemphyxmini0630-pnjunctionsetup}
  A position lies strictly outside the two depleted slabs
\end{definition}
\begin{proof}
  This is the declared model definition of Outside Depletion Region.
\end{proof}
\begin{definition}[Matches Problem Data]
  \label{def:physics:phyx-mini-0630:phyxminiproblems-problemphyxmini0630-matchesproblemdata}
  \lean{PhyXMiniProblems.ProblemPhyXMini0630.MatchesProblemData}
  \uses{def:physics:phyx-mini-0630:phyxminiproblems-problemphyxmini0630-lengthinnanometers, def:physics:phyx-mini-0630:phyxminiproblems-problemphyxmini0630-signedpositioninnanometers, def:physics:phyx-mini-0630:phyxminiproblems-problemphyxmini0630-numberdensitypercubiccentimeter, def:physics:phyx-mini-0630:phyxminiproblems-problemphyxmini0630-pnjunctionsetup}
  Material identities, dopants, orientation, numerical inputs, and the geometry encoded by the labels -a\_p, 0, a\_n, and W. No numerical p-side depth occurs in this interface
\end{definition}
\begin{proof}
  This is the declared model definition of Matches Problem Data.
\end{proof}
\begin{definition}[Matches Primary Figure]
  \label{def:physics:phyx-mini-0630:phyxminiproblems-problemphyxmini0630-matchesprimaryfigure}
  \lean{PhyXMiniProblems.ProblemPhyXMini0630.MatchesPrimaryFigure}
  \uses{def:physics:phyx-mini-0630:phyxminiproblems-problemphyxmini0630-pnjunctionsetup}
  Exact labels, signs, and left-to-right ordering visible in 630.png
\end{definition}
\begin{proof}
  This is the declared model definition of Matches Primary Figure.
\end{proof}
\begin{definition}[Has Physical Parameters]
  \label{def:physics:phyx-mini-0630:phyxminiproblems-problemphyxmini0630-hasphysicalparameters}
  \lean{PhyXMiniProblems.ProblemPhyXMini0630.HasPhysicalParameters}
  \uses{def:physics:phyx-mini-0630:phyxminiproblems-problemphyxmini0630-chargemagnitudeincoulombs, def:physics:phyx-mini-0630:phyxminiproblems-problemphyxmini0630-lengthinnanometers, def:physics:phyx-mini-0630:phyxminiproblems-problemphyxmini0630-numberdensitypercubiccentimeter, def:physics:phyx-mini-0630:phyxminiproblems-problemphyxmini0630-pnjunctionsetup}
  Positivity conditions selecting a nondegenerate physical junction
\end{definition}
\begin{proof}
  This is the declared model definition of Has Physical Parameters.
\end{proof}
\begin{definition}[Satisfies Depletion Region Physics]
  \label{def:physics:phyx-mini-0630:phyxminiproblems-problemphyxmini0630-satisfiesdepletionregionphysics}
  \lean{PhyXMiniProblems.ProblemPhyXMini0630.SatisfiesDepletionRegionPhysics}
  \uses{def:physics:phyx-mini-0630:phyxminiproblems-problemphyxmini0630-chargemagnitudeincoulombs, def:physics:phyx-mini-0630:phyxminiproblems-problemphyxmini0630-axialelectricfieldinvoltspermeter, def:physics:phyx-mini-0630:phyxminiproblems-problemphyxmini0630-lengthincentimeters, def:physics:phyx-mini-0630:phyxminiproblems-problemphyxmini0630-numberdensitypercubiccentimeter, def:physics:phyx-mini-0630:phyxminiproblems-problemphyxmini0630-chargedensityincoulombspercubiccentimeter, def:physics:phyx-mini-0630:phyxminiproblems-problemphyxmini0630-pnjunctionsetup, def:physics:phyx-mini-0630:phyxminiproblems-problemphyxmini0630-inpsidedepletion, def:physics:phyx-mini-0630:phyxminiproblems-problemphyxmini0630-innsidedepletion, def:physics:phyx-mini-0630:phyxminiproblems-problemphyxmini0630-indepletionregion, def:physics:phyx-mini-0630:phyxminiproblems-problemphyxmini0630-outsidedepletionregion}
  The depletion approximation and equilibrium neutrality law. The two fixed charge-density formulas express rho\_p = -e N\_A and rho\_n = e N\_D in consistent centimetre-based units. The areal-neutrality relation N\_A a\_p = N\_D a\_n is a governing equilibrium law relating independent physical inputs; it does not supply the requested numerical value of a\_p
\end{definition}
\begin{proof}
  This is the declared model definition of Satisfies Depletion Region Physics.
\end{proof}
\begin{definition}[Answer Choice]
  \label{def:physics:phyx-mini-0630:phyxminiproblems-problemphyxmini0630-answerchoice}
  \lean{PhyXMiniProblems.ProblemPhyXMini0630.AnswerChoice}
  Labels of the four depletion-depth choices printed with the problem
\end{definition}
\begin{proof}
  This is the declared model definition of Answer Choice.
\end{proof}
\begin{definition}[depth In Nanometers]
  \label{def:physics:phyx-mini-0630:phyxminiproblems-problemphyxmini0630-answerchoice-depthinnanometers}
  \lean{PhyXMiniProblems.ProblemPhyXMini0630.AnswerChoice.depthInNanometers}
  \uses{def:physics:phyx-mini-0630:phyxminiproblems-problemphyxmini0630-answerchoice}
  Nanometre value printed beside each answer label
\end{definition}
\begin{proof}
  This is the declared model definition of depth In Nanometers.
\end{proof}
\begin{definition}[Matches Answer Choice]
  \label{def:physics:phyx-mini-0630:phyxminiproblems-problemphyxmini0630-matchesanswerchoice}
  \lean{PhyXMiniProblems.ProblemPhyXMini0630.MatchesAnswerChoice}
  \uses{def:physics:phyx-mini-0630:phyxminiproblems-problemphyxmini0630-lengthinnanometers, def:physics:phyx-mini-0630:phyxminiproblems-problemphyxmini0630-pnjunctionsetup, def:physics:phyx-mini-0630:phyxminiproblems-problemphyxmini0630-answerchoice}
  A displayed choice matches the independently modeled p-side depth
\end{definition}
\begin{proof}
  This is the declared model definition of Matches Answer Choice.
\end{proof}
% --- Archon declaration topology end ---
% --- Archon physics formalization source end ---
